\par\Needspace{75mm}
\tagstructbegin{tag=Figure,alttext={视觉到语言的两步桥接。四行三维视觉特征乘三乘四桥接矩阵，得到四行四维视觉前缀。再在行方向拼接两行四维文字向量，形成六行四维输入。每行宽度都为四，前四行来自图片，后两行来自文字。}}
\begin{center}
\begin{tikzpicture}[>=Latex,x=1mm,y=1mm,
 block/.style={draw=linegray,fill=paper,rounded corners=2pt,align=center,font=\small,minimum height=15mm,text width=35mm}]
\node[block,fill=linegray!15] (image) at (0,0) {四行视觉特征\\$H_{\rm img}:4\times3$};
\node[block] (bridge) at (53,0) {桥接矩阵\\$W_{\rm bridge}:3\times4$};
\node[block,fill=linegray!15] (prefix) at (106,0) {四行视觉前缀\\$H_{\rm prefix}:4\times4$};
\draw[->,thick,flowcolor] (image)--(bridge);
\draw[->,thick,flowcolor] (bridge)--(prefix);
\node[block,fill=paper] (text) at (0,-36) {两行文字向量\\$E_{\rm text}:2\times4$};
\node[block,text width=48mm,minimum height=30mm] (input) at (87,-45) {位置 1--4：视觉前缀\\[2mm]\textcolor{flowcolor}{位置 5--6：文字向量}\\[2mm]$X_{\rm LLM}:6\times4$};
\draw[->,thick,flowcolor] (prefix.south)--(input.north -| prefix.south);
\draw[->,thick,flowcolor] (text.east)--(input.west |- text.east);
\node[font=\footnotesize,text=muted,align=center] at (16,-62) {先对齐列宽，再增加行数\\最后一个文字位置可读取六行};
\end{tikzpicture}
\captionof{figure}{视觉前缀接入语言模型：桥接改变特征宽度，拼接增加序列长度}
\label{fig:vision-language-prefix}
\end{center}
\tagstructend
